\documentclass[11pt]{article}
% Shared preamble for the printable mini lessons (1.4 onward).
%   \documentclass[11pt]{article}
%   % Shared preamble for the printable mini lessons (1.4 onward).
%   \documentclass[11pt]{article}
%   % Shared preamble for the printable mini lessons (1.4 onward).
%   \documentclass[11pt]{article}
%   \input{../preamble}
% Build the key with:  pdflatex -jobname=<name>_KEY "\def\iskey{}\input{<name>}"
% Figures shared with the web page live in figures/Unit_X/<name>.tex and are
% pulled in with \fig{Unit_X/<name>}; the standalone package drops their own
% preamble, so their macros come from figures/preamble.tex, loaded here.
\input{../../figures/preamble}
\usepackage[margin=0.75in]{geometry}
\usepackage{array}
\usepackage{enumitem}
\usepackage{amsmath}
\usepackage[helvet]{sfmath}  % math in the same sans font as the text
\usepackage{standalone}
\pagestyle{empty}
\setlength{\parindent}{0pt}
\setlist[enumerate,1]{leftmargin=1.5em,itemsep=1em}
\setlist[enumerate,2]{label=(\alph*),leftmargin=2em,itemsep=2.2em}

\newif\ifkey
\ifdefined\iskey \keytrue \fi
\colorlet{keyred}{red!75!black}
\newcommand{\ans}[1]{\ifkey\leavevmode{\color{keyred}#1}\fi}
% vertical work space on the student copy, less on the key
\newcommand{\work}[1]{\ifkey\vspace{0.6em}\else\vspace{#1}\fi}
\newcommand{\fig}[1]{\input{../../figures/#1}}
% Title line: \lessontitle{Quartiles, IQR, and Box Plots}
\newcommand{\lessontitle}[1]{{\Large\bfseries #1\ifkey\ \textcolor{keyred}{KEY}\fi}\hfill Name \underline{\hspace{2in}}\par\vspace{0.6em}}
\newcommand{\partner}{\par\vspace{0.4em}Work with a partner. Compare answers when you finish each problem.\par}

% Build the key with:  pdflatex -jobname=<name>_KEY "\def\iskey{}\input{<name>}"
% Figures shared with the web page live in figures/Unit_X/<name>.tex and are
% pulled in with \fig{Unit_X/<name>}; the standalone package drops their own
% preamble, so their macros come from figures/preamble.tex, loaded here.
% Shared preamble for web figures.  Each figure is a standalone .tex that
% \input's this file; scripts/make-figures.sh turns it into an SVG.
\usepackage[T1]{fontenc}
\usepackage{helvet}
\renewcommand{\familydefault}{\sfdefault}
\usepackage{tikz}
\usetikzlibrary{decorations.pathreplacing}
% axis numbers in the text font, not math italic
\pgfkeys{/pgf/number format/.cd,assume math mode=true,1000 sep={}}

% dot plot: \dotplot{min}{max}{value/count,...}{axis label}
\newcommand{\dotplot}[4]{%
\begin{tikzpicture}[x=0.9cm,y=0.33cm]
  \draw[thick] (#1-0.5,0) -- (#2+0.5,0);
  \foreach \t in {#1,...,#2} {\draw (\t,0.25) -- (\t,-0.25) node[below] {\t};}
  \foreach \v/\c in {#3} {\ifnum\c>0 \foreach \k in {1,...,\c} {\fill (\v,\k*1.15) circle (0.16cm);}\fi}
  \node at ({(#1+#2)/2},-2.3) {#4};
\end{tikzpicture}}

% shape sketch histogram (12 bins): \shapehist{heights}{label}
\newcommand{\shapehist}[2]{%
  \foreach \h [count=\i from 0] in {#1} \draw[fill=gray!30] (\i,0) rectangle (\i+1,\h);
  \draw (-0.3,0) -- (12.3,0);
  \node[font=\bfseries] at (6,-2.6) {#2};}
\newcommand{\meanmark}[1]{\fill (#1,0) -- ++(-0.3,-0.75) -- ++(0.6,0) -- cycle;}
\newcommand{\medmark}[2]{\draw[very thick,dashed] (#1,0) -- (#1,#2);}

% number line with minor ticks: \nline{min}{max}{label step}{axis label}
\newcommand{\nline}[4]{%
  \draw[thick] (#1-0.5,0) -- (#2+0.5,0);
  \foreach \t in {#1,...,#2} \draw (\t,0.12) -- (\t,-0.12);
  \pgfmathtruncatemacro{\nxt}{#1+#3}
  \foreach \t in {#1,\nxt,...,#2} \draw (\t,0.2) -- (\t,-0.2) node[below] {\t};
  \node at ({(#1+#2)/2},-0.95) {#4};}
% box plot above the line: \bplot{min}{Q1}{med}{Q3}{max}
\newcommand{\bplot}[5]{%
  \draw[thick] (#1,1.0) -- (#2,1.0);
  \draw[thick] (#4,1.0) -- (#5,1.0);
  \draw[thick] (#1,0.75) -- (#1,1.25);
  \draw[thick] (#5,0.75) -- (#5,1.25);
  \draw[thick,fill=gray!20] (#2,0.55) rectangle (#4,1.45);
  \draw[very thick] (#3,0.55) -- (#3,1.45);}

% ---- added for 1.4 onward ----
% axis on y=0: \axis{min}{max}{label step}{tick step}{axis label}
\newcommand{\axis}[5]{%
  \draw[thick] (#1,0) -- (#2,0);
  \pgfmathsetmacro{\axT}{#1+#4}\foreach \t in {#1,\axT,...,#2} \draw (\t,0.1) -- (\t,-0.1);
  \pgfmathsetmacro{\axL}{#1+#3}\foreach \t in {#1,\axL,...,#2}
    \draw (\t,0.18) -- (\t,-0.18) node[below] {\pgfmathprintnumber{\t}};
  \node at ({(#1+#2)/2},-0.95) {#5};}
% box plot centred at height y: \bplotat{y}{min}{Q1}{med}{Q3}{max}
\newcommand{\bplotat}[6]{%
  \draw[thick] (#2,#1) -- (#3,#1);
  \draw[thick] (#5,#1) -- (#6,#1);
  \draw[thick] (#2,#1-0.22) -- (#2,#1+0.22);
  \draw[thick] (#6,#1-0.22) -- (#6,#1+0.22);
  \draw[thick,fill=gray!20] (#3,#1-0.42) rectangle (#5,#1+0.42);
  \draw[very thick] (#4,#1-0.42) -- (#4,#1+0.42);}
% an outlier dot: \outlier{x}{y}
\newcommand{\outlier}[2]{\fill (#1,#2) circle (2.4pt);}

\usepackage[margin=0.75in]{geometry}
\usepackage{array}
\usepackage{enumitem}
\usepackage{amsmath}
\usepackage[helvet]{sfmath}  % math in the same sans font as the text
\usepackage{standalone}
\pagestyle{empty}
\setlength{\parindent}{0pt}
\setlist[enumerate,1]{leftmargin=1.5em,itemsep=1em}
\setlist[enumerate,2]{label=(\alph*),leftmargin=2em,itemsep=2.2em}

\newif\ifkey
\ifdefined\iskey \keytrue \fi
\colorlet{keyred}{red!75!black}
\newcommand{\ans}[1]{\ifkey\leavevmode{\color{keyred}#1}\fi}
% vertical work space on the student copy, less on the key
\newcommand{\work}[1]{\ifkey\vspace{0.6em}\else\vspace{#1}\fi}
\newcommand{\fig}[1]{\input{../../figures/#1}}
% Title line: \lessontitle{Quartiles, IQR, and Box Plots}
\newcommand{\lessontitle}[1]{{\Large\bfseries #1\ifkey\ \textcolor{keyred}{KEY}\fi}\hfill Name \underline{\hspace{2in}}\par\vspace{0.6em}}
\newcommand{\partner}{\par\vspace{0.4em}Work with a partner. Compare answers when you finish each problem.\par}

% Build the key with:  pdflatex -jobname=<name>_KEY "\def\iskey{}\input{<name>}"
% Figures shared with the web page live in figures/Unit_X/<name>.tex and are
% pulled in with \fig{Unit_X/<name>}; the standalone package drops their own
% preamble, so their macros come from figures/preamble.tex, loaded here.
% Shared preamble for web figures.  Each figure is a standalone .tex that
% \input's this file; scripts/make-figures.sh turns it into an SVG.
\usepackage[T1]{fontenc}
\usepackage{helvet}
\renewcommand{\familydefault}{\sfdefault}
\usepackage{tikz}
\usetikzlibrary{decorations.pathreplacing}
% axis numbers in the text font, not math italic
\pgfkeys{/pgf/number format/.cd,assume math mode=true,1000 sep={}}

% dot plot: \dotplot{min}{max}{value/count,...}{axis label}
\newcommand{\dotplot}[4]{%
\begin{tikzpicture}[x=0.9cm,y=0.33cm]
  \draw[thick] (#1-0.5,0) -- (#2+0.5,0);
  \foreach \t in {#1,...,#2} {\draw (\t,0.25) -- (\t,-0.25) node[below] {\t};}
  \foreach \v/\c in {#3} {\ifnum\c>0 \foreach \k in {1,...,\c} {\fill (\v,\k*1.15) circle (0.16cm);}\fi}
  \node at ({(#1+#2)/2},-2.3) {#4};
\end{tikzpicture}}

% shape sketch histogram (12 bins): \shapehist{heights}{label}
\newcommand{\shapehist}[2]{%
  \foreach \h [count=\i from 0] in {#1} \draw[fill=gray!30] (\i,0) rectangle (\i+1,\h);
  \draw (-0.3,0) -- (12.3,0);
  \node[font=\bfseries] at (6,-2.6) {#2};}
\newcommand{\meanmark}[1]{\fill (#1,0) -- ++(-0.3,-0.75) -- ++(0.6,0) -- cycle;}
\newcommand{\medmark}[2]{\draw[very thick,dashed] (#1,0) -- (#1,#2);}

% number line with minor ticks: \nline{min}{max}{label step}{axis label}
\newcommand{\nline}[4]{%
  \draw[thick] (#1-0.5,0) -- (#2+0.5,0);
  \foreach \t in {#1,...,#2} \draw (\t,0.12) -- (\t,-0.12);
  \pgfmathtruncatemacro{\nxt}{#1+#3}
  \foreach \t in {#1,\nxt,...,#2} \draw (\t,0.2) -- (\t,-0.2) node[below] {\t};
  \node at ({(#1+#2)/2},-0.95) {#4};}
% box plot above the line: \bplot{min}{Q1}{med}{Q3}{max}
\newcommand{\bplot}[5]{%
  \draw[thick] (#1,1.0) -- (#2,1.0);
  \draw[thick] (#4,1.0) -- (#5,1.0);
  \draw[thick] (#1,0.75) -- (#1,1.25);
  \draw[thick] (#5,0.75) -- (#5,1.25);
  \draw[thick,fill=gray!20] (#2,0.55) rectangle (#4,1.45);
  \draw[very thick] (#3,0.55) -- (#3,1.45);}

% ---- added for 1.4 onward ----
% axis on y=0: \axis{min}{max}{label step}{tick step}{axis label}
\newcommand{\axis}[5]{%
  \draw[thick] (#1,0) -- (#2,0);
  \pgfmathsetmacro{\axT}{#1+#4}\foreach \t in {#1,\axT,...,#2} \draw (\t,0.1) -- (\t,-0.1);
  \pgfmathsetmacro{\axL}{#1+#3}\foreach \t in {#1,\axL,...,#2}
    \draw (\t,0.18) -- (\t,-0.18) node[below] {\pgfmathprintnumber{\t}};
  \node at ({(#1+#2)/2},-0.95) {#5};}
% box plot centred at height y: \bplotat{y}{min}{Q1}{med}{Q3}{max}
\newcommand{\bplotat}[6]{%
  \draw[thick] (#2,#1) -- (#3,#1);
  \draw[thick] (#5,#1) -- (#6,#1);
  \draw[thick] (#2,#1-0.22) -- (#2,#1+0.22);
  \draw[thick] (#6,#1-0.22) -- (#6,#1+0.22);
  \draw[thick,fill=gray!20] (#3,#1-0.42) rectangle (#5,#1+0.42);
  \draw[very thick] (#4,#1-0.42) -- (#4,#1+0.42);}
% an outlier dot: \outlier{x}{y}
\newcommand{\outlier}[2]{\fill (#1,#2) circle (2.4pt);}

\usepackage[margin=0.75in]{geometry}
\usepackage{array}
\usepackage{enumitem}
\usepackage{amsmath}
\usepackage[helvet]{sfmath}  % math in the same sans font as the text
\usepackage{standalone}
\pagestyle{empty}
\setlength{\parindent}{0pt}
\setlist[enumerate,1]{leftmargin=1.5em,itemsep=1em}
\setlist[enumerate,2]{label=(\alph*),leftmargin=2em,itemsep=2.2em}

\newif\ifkey
\ifdefined\iskey \keytrue \fi
\colorlet{keyred}{red!75!black}
\newcommand{\ans}[1]{\ifkey\leavevmode{\color{keyred}#1}\fi}
% vertical work space on the student copy, less on the key
\newcommand{\work}[1]{\ifkey\vspace{0.6em}\else\vspace{#1}\fi}
\newcommand{\fig}[1]{\input{../../figures/#1}}
% Title line: \lessontitle{Quartiles, IQR, and Box Plots}
\newcommand{\lessontitle}[1]{{\Large\bfseries #1\ifkey\ \textcolor{keyred}{KEY}\fi}\hfill Name \underline{\hspace{2in}}\par\vspace{0.6em}}
\newcommand{\partner}{\par\vspace{0.4em}Work with a partner. Compare answers when you finish each problem.\par}

\newcommand{\blank}[1]{\ans{#1}}

\begin{document}

\lessontitle{Two-Way Tables and Relative Frequency}

60 students answered two questions: \textit{Do you play a sport?} \textit{Do you have a job?}
A \textbf{two-way table} shows the answers to both questions at once.

\vspace{0.6em}
\begin{minipage}[c]{0.5\linewidth}
\renewcommand{\arraystretch}{1.4}
\begin{tabular}{|l|c|c|c|}\hline
 & \textbf{Job} & \textbf{No job} & \textbf{Total} \\ \hline
\textbf{Sport} & 8 & 22 & 30 \\ \hline
\textbf{No sport} & 12 & 18 & 30 \\ \hline
\textbf{Total} & 20 & 40 & 60 \\ \hline
\end{tabular}
\end{minipage}\hfill
\begin{minipage}[c]{0.45\linewidth}\centering
\scalebox{0.75}{\fig{Unit_3/3.1_venn_sport_job}}
\end{minipage}

\vspace{0.6em}
\begin{tabular}{@{}ll}
Each \textbf{cell} counts students who match \textbf{both} labels. & 8 students play a sport \textbf{and} have a job.\\
The \textbf{totals} are in the margins. & 30 students play a sport. 20 have a job.\\
The \textbf{grand total} is the corner. & 60 students in all.\\
\end{tabular}

\vspace{0.4em}
\textbf{Check:} Each row adds to its total. Each column adds to its total. \quad $8 + 22 = 30$, \; $8 + 12 = 20$.

\vspace{0.4em}
\textbf{Venn diagram:} the overlap is ``both'' (8). Outside the circles is ``neither'' (18).

\vspace{1em}
\textbf{Relative frequency table.} Divide every count by the \textbf{grand total}. This gives the percent of \textbf{all} students.

\vspace{0.4em}
\renewcommand{\arraystretch}{2.3}
\begin{tabular}{|l|c|c|c|}\hline
 & \textbf{Job} & \textbf{No job} & \textbf{Total} \\ \hline
\textbf{Sport} & $\dfrac{8}{60} \approx 13\%$ & $\dfrac{22}{60} \approx 37\%$ & $\dfrac{30}{60} = 50\%$ \\ \hline
\textbf{No sport} & $\dfrac{12}{60} = 20\%$ & $\dfrac{18}{60} = 30\%$ & $\dfrac{30}{60} = 50\%$ \\ \hline
\textbf{Total} & $\dfrac{20}{60} \approx 33\%$ & $\dfrac{40}{60} \approx 67\%$ & 100\% \\ \hline
\end{tabular}

\vspace{0.4em}
About 13\% of \textbf{all} the students play a sport and have a job. The four inside cells add to 100\%.

\partner

% ---------------- PROBLEMS ----------------
\newpage
\begin{enumerate}

\item 80 students said how they get to school.

\begin{center}
\renewcommand{\arraystretch}{1.4}
\begin{tabular}{|l|c|c|c|c|}\hline
 & \textbf{Bus} & \textbf{Car} & \textbf{Walk} & \textbf{Total} \\ \hline
\textbf{9th grade} & 18 & 12 & 10 & 40 \\ \hline
\textbf{10th grade} & 14 & 20 & 6 & 40 \\ \hline
\textbf{Total} & 32 & 32 & 16 & 80 \\ \hline
\end{tabular}
\end{center}
\begin{enumerate}[itemsep=1.4em]
  \item How many 10th graders walk? \ans{\quad 6}
  \item How many students ride the bus? \ans{\quad 32}
  \item How many students are 9th graders who come by car? \ans{\quad 12}
  \item How many students do \textbf{not} walk? \ans{\quad $80 - 16 = 64$}
\end{enumerate}
\work{1.5em}

\item Fill in the missing numbers. Students said if they ate breakfast today.

\begin{center}
\renewcommand{\arraystretch}{1.6}
\begin{tabular}{|l|p{0.8in}|p{0.8in}|p{0.8in}|}\hline
 & \textbf{Yes} & \textbf{No} & \textbf{Total} \\ \hline
\textbf{11th grade} & 15 & \blank{20} & 35 \\ \hline
\textbf{12th grade} & \blank{12} & 10 & \blank{22} \\ \hline
\textbf{Total} & 27 & \blank{30} & 57 \\ \hline
\end{tabular}
\end{center}
\work{2em}

\item 50 people were asked about pets. 30 like cats. 26 like dogs. 12 like both.
\begin{enumerate}[itemsep=1em]
  \item Fill in the two-way table.

  \begin{center}
  \renewcommand{\arraystretch}{1.6}
  \begin{tabular}{|l|p{0.9in}|p{0.9in}|p{0.9in}|}\hline
   & \textbf{Likes dogs} & \textbf{No dogs} & \textbf{Total} \\ \hline
  \textbf{Likes cats} & \blank{12} & \blank{18} & \blank{30} \\ \hline
  \textbf{No cats} & \blank{14} & \blank{6} & \blank{20} \\ \hline
  \textbf{Total} & \blank{26} & \blank{24} & \blank{50} \\ \hline
  \end{tabular}
  \end{center}
  \item How many people like neither cats nor dogs? \ans{\quad 6}
\end{enumerate}

\newpage
\item This Venn diagram shows which sports 50 students play.
\begin{center}\scalebox{0.75}{\fig{Unit_3/3.1_venn_sports}}\end{center}
Make a two-way table from the Venn diagram.
\begin{center}
\renewcommand{\arraystretch}{1.6}
\begin{tabular}{|l|p{1in}|p{1in}|p{1in}|}\hline
 & \textbf{Soccer} & \textbf{No soccer} & \textbf{Total} \\ \hline
\textbf{Basketball} & \blank{5} & \blank{15} & \blank{20} \\ \hline
\textbf{No basketball} & \blank{10} & \blank{20} & \blank{30} \\ \hline
\textbf{Total} & \blank{15} & \blank{35} & \blank{50} \\ \hline
\end{tabular}
\end{center}
\work{1.5em}

\item Make a relative frequency table for problem 1. Divide each count by 80. Write percents.
\begin{center}
\renewcommand{\arraystretch}{1.6}
\begin{tabular}{|l|p{0.8in}|p{0.8in}|p{0.8in}|p{0.8in}|}\hline
 & \textbf{Bus} & \textbf{Car} & \textbf{Walk} & \textbf{Total} \\ \hline
\textbf{9th grade} & \blank{22.5\%} & \blank{15\%} & \blank{12.5\%} & \blank{50\%} \\ \hline
\textbf{10th grade} & \blank{17.5\%} & \blank{25\%} & \blank{7.5\%} & \blank{50\%} \\ \hline
\textbf{Total} & \blank{40\%} & \blank{40\%} & \blank{20\%} & \blank{100\%} \\ \hline
\end{tabular}
\end{center}
\begin{enumerate}[itemsep=1.4em]
  \item What percent of all the students are 10th graders who come by car? \ans{\quad 25\%}
  \item What percent of all the students walk? \ans{\quad 20\%}
\end{enumerate}

\end{enumerate}

\end{document}
